In this section, we will recall definitions associated with cone
deformations and review relevant results from work of Hodgson and
Kerckhoff (\cite{hk:cone-rigid}, \cite{hk:univ}).

\begin{define}
Let $M$ be a 3--manifold containing a closed 1--manifold $\Sigma$.  A
\emph{hyperbolic cone structure} on $M$ is an incomplete hyperbolic
structure on $X=M-\Sigma$, whose completion is singular along
$\Sigma$.  A cross section of a component of $\Sigma$ is a hyperbolic
cone with angle $\alpha$, where $\alpha$ is constant along that
component.  (For a more complete description of the metric on $M$, see
\cite{hk:cone-rigid}.)  The set $\Sigma$ is called the \emph{singular
locus} of $M$.  $M$ with this structure is called a \emph{hyperbolic
cone manifold}.  A \emph{hyperbolic cone deformation} on $M$ is a
smooth one--parameter family of hyperbolic cone structures $X_t$.
\label{def:cone-deform}
\end{define}

\subsection{Infinitesimal deformations}
Now, we may study deformations by considering possible ``derivatives
at time $t$'' of $X_t$, or infinitesimal deformations of hyperbolic
structure,
%% Aug 6, 2007.  Emphasize time dependency.
for any time $t$.
%%
Each of these infinitesimal deformations corresponds to an element in
a certain cohomology group $H^1(X; E)$, where recall $X=M-\Sigma$.  We
will be manipulating elements of $H^1(X;E)$, so we recall some
information on their structure.


%% April 7.  This subsection of cone deformations section gives
%% description of $\omega$ and $\eta$ and their decompositions in
%% tubular and horoball neighborhoods.


\subsubsection{Cohomology}

Given $\omega$ in $H^1(X;E)$, $\omega$ is a one--form on $X$ with
values in the vector bundle $E$ of infinitesimal isometries of $X$.
The bundle $E$ has fiber $sl(2,\CC)$, which is the {L}ie algebra of
infinitesimal isometries of $\HH^3$.  Recall that this {L}ie algebra
has a complex structure.  Geometrically, if $s$ represents an
infinitesimal translation in the direction of $s$, then $i s$
represents an infinitesimal rotation with axis in the direction of
$s$.  Thus on $X$ we can identify $E$ with the complexified tangent
bundle $TX \otimes \CC$, and write any $\omega \in H^1(X;E)$ in
complex form: $\omega = v + iw$.

By the {H}odge theorem proved in \cite{hk:cone-rigid}, in each
cohomology class of $H^1(X;E)$ there is a \emph{harmonic}
representative of the form
$$\omega = \eta + i*D\eta,$$
where $\eta$ is a unique $TX$-valued
1--form on $X = M-\Sigma$ that satisfies:
$$D^*\eta=0,$$
\begin{equation}
D^*D\eta + \eta =0.
\label{eqn:d*deta}
\end{equation}
Here $D$ is the exterior covariant derivative on such forms and $D^*$
is its adjoint.

%% Aug 6, 2007.  Emphasize time dependency.
Thus we may assume that our cone deformation is chosen such that at
each time $t$, the infinitesimal deformation of hyperbolic structure
of $X_t$ corresponds to a \emph{harmonic} element $\omega$ in
$H^1(X;E)$.
%%

\subsubsection{Tubular and Horoball neighborhoods}

For purposes of this paper, we are particularly interested in the
effect of the deformation in a neighborhood of the singular locus or
in a horoball neighborhood of a cusp.  So in this subsection, we will
review a decomposition of $\omega$ (and $\eta$) into a nice form in
these types of neighborhoods.

Select a component of the singular locus $\Sigma$ of $M$.  Let $V$ be
a tubular neighborhood of that component, with tube radius $R$.  $V$
is a solid torus with a singular core curve.  During the cone
deformation, the geometric structure on $V$ is changing in meridional,
longitudinal, and radial directions.  We can write $\omega$ in $V$ to
reflect this:
\begin{equation}
\omega = \omega_0 + \omega_c.
\label{eqn:omega-decomp}
\end{equation}
Here only $\omega_0 = \eta_0 + i*D\eta_0$
changes the holonomy of the meridian and longitude on the torus
$\partial V$, and $\omega_c$ is a correction term.

We can further decompose $\omega_0$ as follows:
\begin{equation}
\omega_0 = u\, \omega_m + (x+iy)\,\omega_\ell.
\label{eqn:omega_0-decomp}
\end{equation}
Here, $u$, $x$, and $y$ are real numbers.  The 1--forms $\omega_m$ and
$\omega_\ell$ are standard forms, calculated in \cite{hk:cone-rigid},
which depend only on the tube radius $R$ of $V$.  The form $\omega_m =
\eta_m + i*D\eta_m$ gives the change in the meridional direction.  In
particular, it represents the infinitesimal deformation which
decreases the cone angle but doesn't change the real part of the
complex length of the meridian.  The form $\omega_\ell = \eta_\ell
+i*D\eta_\ell$ stretches the singular locus, but leaves the holonomy
of the meridian unchanged.

Since we will use the explicit formula for the standard form
$\omega_\ell$ in a calculation in a later section, we will restate
that form here.  Let $e_1$, $e_2$, and $e_3$ be an orthonormal frame
for $V$ in cylindrical coordinates, with $e_1$ pointing in the radial
direction and $e_2$ tangent to the meridian.  Let $R$ be the radius of
the tube $V$.  Recall $\omega_\ell$ is a ($TX\otimes\CC$)--valued
1--form, which can be viewed as an element of $Hom(TX, TX\otimes\CC)$.
Thus it can be described as a matrix in the cylindrical coordinates
$e_1$, $e_2$, $e_3$.
\begin{equation}
  \omega_\ell(R) = \left[
    \begin{array}{ccccc}
      (\cosh R)^{-2} & \hspace{.1in} & 0 &\hspace{.1in}& 0 \\
      0 && -1 && -i\tanh R \\
      0 && -i\tanh R && 1+ (\cosh R)^{-2}
    \end{array}
    \right]
  \label{eqn:omega_l}
\end{equation}


Now, so far, our decomposition of $\omega$ has been in a tubular
neighborhood $V$ of a component of the singular locus.  We can do a
similar decomposition of $\omega$ in a horoball neighborhood $U$ of a
cusp.  In fact, we may consider $U$ as a tubular neighborhood of a
1--dimensional submanifold of $M$ on which the cone angle remains $0$
throughout the deformation, and the tube radius is infinite.  Again we
may decompose $\omega$ as $\omega_0 + \omega_c$, where only $\omega_c$
affects the holonomy of the torus $\partial U$.  But now, since the
cone angle at the core of $U$ is not changing, the meridional part of
$\omega_0$ is identically $0$.  Thus in this case,
$$\omega_0 = (x + iy)\, \omega_\ell,$$
where again $x$ and $y$ are real numbers, and $\omega_\ell$ is a
standard form.  Since the tube radius of $U$ is infinite, to write
down the formula for $\omega_\ell$ explicitly we let $R$ go to
infinity in equation (\ref{eqn:omega_l}).
\begin{equation}
  \omega_\ell(\infty) = \left[
    \begin{array}{ccc}
      0 & 0 & 0\\
      0 & -1 & -i\\
      0 & -i & 1
    \end{array}\right]
  \label{eqn:omega_l-cusp}
\end{equation}

%%%%%%%%%%%%%%%%%%%%%%%%%%%%%%%%%%%%%%%%%%%%%%%%%%%%%%%%%%%%%%%%%


\subsection{Boundary terms}
\label{subsec:boundary-terms}

Our proof of Theorem \ref{thm:main-cusp} will involve certain boundary
terms.  In this section, we define these terms.  We explain how they
arise and recall certain properties.  They will be used in the
proof of Theorem \ref{thm:main-cusp}.


%% April 7:  In this subsection, we define boundary terms and present
%% inequalities relating them.

Using results of the last section, we may always assume that we have a
harmonic representative $\omega = \eta + i*D\eta$.  Thus $\eta$
satisfies equation (\ref{eqn:d*deta}).

Let $N$ be any submanifold of $X = M-\Sigma$ with boundary $\partial
N$ oriented by the outward normal.  When we take the $L^2$ inner
product of equation (\ref{eqn:d*deta}) with $\eta$ and integrate by
parts over $N$, we obtain the boundary formula (\cite[Lemma
2.3]{hk:univ}):
\begin{equation}
  B(N)= \int_{\partial N} \eta \wedge *D\eta = \| \eta \|^2_N + \|
  *D\eta \|^2_N 
  \label{eqn:BN}
\end{equation}
Here $\| \cdot \|$ denotes the $L^2$--norm on $N$.  We will apply
equation (\ref{eqn:BN}) to submanifolds of $X=M-\Sigma$ involving
tubular or horoball neighborhoods.

First, we will introduce some notation.  Let $V_j$ be a tubular
neighborhood of a component (the $j$-th component) of the singular
locus, or a horoball neighborhood of a cusp.  Using similar notation
to that of \cite{hk:univ}, define
$$b_{V_j}(\zeta,\xi)=\int_{\partial V_j}*D\zeta\wedge \xi,$$
where
$\partial V_j$ is oriented such that the tangent vector orthogonal to
$\partial V_j$ is an outward normal when viewed from $V_j$.  We will
be considering $b_{V_j}(\eta, \eta)$ in this paper, and we review
results concerning this term.

The decomposition (\ref{eqn:omega-decomp}) of $\omega$ in $V_j$ into
$\omega_0 + \omega_c$, and $\eta$ into $\eta_0 + \eta_c$ decomposes
the term $b_{V_j}(\eta, \eta)$:
\begin{equation}
b_{V_j}(\eta, \eta) = b_{V_j}(\eta_0, \eta_0) +
b_{V_j}(\eta_c, \eta_c).
\label{eqn:bV-decomp}
\end{equation}
That is, cross terms vanish (\cite[Lemma 2.5]{hk:univ}).
Additionally, the term $b_{V_j}(\eta_c, \eta_c)$ is actually
non-positive (\cite[Lemma 2.6]{hk:univ}), so we find
\begin{equation}
  b_{V_j}(\eta, \eta) \leq b_{V_j}(\eta_0, \eta_0)
\label{eqn:eta<eta_0}
\end{equation}


Now, let $V$ be a tubular neighborhood of the entire singular locus
$\Sigma$, as well as horoball neighborhoods of cusps.  Thus $V$ is a
union of tubular neighborhoods $V_j$ of components $\Sigma_j$ of the
singular locus and horoball neighborhoods $U_j$.  Then letting $N=X-V$
in equation (\ref{eqn:BN}),
\begin{eqnarray}
B(X-V)  & = &\| \eta \|^2_{X-V} + \|*D\eta\|^2_{X-V} =
 \int_{\partial V} *D\eta \wedge \eta \\
 & = &
 \sum_j b_{V_j} (\eta, \eta) + \sum_k b_{U_j}(\eta,\eta).
\label{eqn:sumbj=eta}
\end{eqnarray}
The first sum is over all components of the singular locus, the second
over all horoball neighborhoods of cusps.  This equation, along with
equation (\ref{eqn:eta<eta_0}), implies
\begin{equation}
  0 \leq  \sum_j b_{V_j} (\eta, \eta)+ \sum_k
  b_{U_k}(\eta,\eta)  \leq \sum_j b_{V_j} (\eta_0,
  \eta_0) + \sum_k b_{U_k} (\eta_0, \eta_0).
\label{eqn:sum>0}
\end{equation}

Notice $\eta_0 = (\eta_0)_j$ depends on the neighborhood $V_j$ or
$U_j$.  However, using the notation $b_{V_j}(\eta_0, \eta_0)$, it
should be clear that we are referring to the decomposition of $\eta$
particular to $V_j$ in this context, and similarly for $U_j$.  We will
further simplify notation, using the following definition.

\begin{define}
  Let $\Sigma_j$ be a component of the singular locus $\Sigma$.  Let
  $V_j$ be a tubular neighborhood of $\Sigma_j$ of radius $R$.  We let
  $b_j$ be the boundary term:
  $$b_j = b_{V_j}(\eta_0, \eta_0).$$
  \label{def:bj}
\end{define}

\begin{remark}
Note that for our applications of the results of this section, the
one--form $\omega$ corresponds to an infinitesimal deformation of
hyperbolic structure at time $t$.  Thus $\omega$, $\eta$, $\eta_0$,
$b_j$, etc. will all depend on time $t$ in our applications in future
sections.

To avoid a notational nightmare, in the sequel we will assume the time
$t$ has been fixed, and continue writing these terms without reference
to $t$, except occasionally where we feel recalling the dependency on
$t$ will help avoid confusion.
\end{remark}

